\section{Climate, Oxygen, and the Future of Marine Biodiversity}

\begin{itemize}
	\item The ocean enabled the diversification of life on Earth by adding
		O$_2$ to the atmosphere.
	\item Human industrialization is intensifying the aerobic challenges
		to marine ecosystems.
	\item Historical observations report an 2\% decline in upper-ocean
		O$_2$.
	\item The ocean retains under 2\% of global O$_2$ inventory, while
		98\% stays in the atmosphere.
	\item The constraining role of O$_2$ for ocean species arises from
		two kinematic factors: lower solubility of gases in water and
		slower circulation of the ocean than of atmosphere.
	\item Industrialization of human societies produced two negative
		consequences for the ocean's O$_2$ reservoir: (1) nutrients
		mobilized for agriculture are carried to coastal oceans,
		stimulating phyloplankton growth and thus higher O$_2$
		respiratory demand (2) global warming causes O$_2$ depletion
		throughout the ocean
	\item Coastal regions make up only 7\% of ocean's area but contain
		15\% of it's net primary productivity and 90\% of it's animal
		life.
	\item 
\end{itemize}

\subsection{Glossary}

\subsubsection{O$_2$ solubility in seawater (K$_H$)}
O$_2$ in seawater when in equilibrium with the O$_2$ in a gas phase, as a
function of temperature, pressure, and salinity

\subsubsection{O$_2$ partial pressure ($p$O$_2$)}
the ratio of O$_2$ concentration to O$_2$ solubility

\subsubsection{hypoxia}
the condition wherein O$_2$ is insufficient to meet metabolic demand, which
varies with species and temperature, in contrat to fixed definition
commonly used in ocanography (O$_2$ < 60$\mu$M)

\subsubsection{Eutrophication}
increased growth, primary production, and biomass of algae in response to
nutrient enrichment in the environment (e.g. sewage, fertilizer, and
detergents)

\subsubsection{Metabolic storm}
an anomalous ocean weather event in which high temperatures cause elevated
organismal O$_2$ demand simultaneous with decline in O$_2$ available in
seawater

\subsubsection{Temperature-dependent hypoxia}
the threshold for hypoxia and its variation with temperature

\subsubsection{Oxygen minimum zone (OMZ)}
a vertical minimum in a water column; OMZs are globally prevalent but are
strongest in the tropical thermoclines of the Indian and Pacific basins, where
O$_2$ reaches traditionally defined hypoxic concentrations

\subsubsection{Suboxia}
O$_2$ < 5 $\mu$M, corresponding to the appearance of nitrogen deficits
relative to phosphorous in seawater and the activation of anaerobic microbial
metabolisms; used interchangeably with anoxia

\subsubsection{Apparent O$_2$ utilization (AOU)}
the amount of O$_2$ consumed by respiration over the transport time of a
water parcel from the surface ocean to the interior

\subsubsection{Critical O$_2$ pressure in seawater}
the ambient $p$O$_2$ when an organism's O$_2$ supply balances its resting
metabolic O$_2$ demand
